\documentclass[12pt,a4paper]{article}
\usepackage[margin=2.5cm]{geometry}
\usepackage[T1]{fontenc}
\usepackage[utf8]{inputenc}
\usepackage{microtype}
\usepackage[colorlinks=true,urlcolor=blue]{hyperref}
\setlength{\parindent}{0pt}
\setlength{\parskip}{0.8em}

\begin{document}

3 May 2026

Editorial Office\\
\textit{Scientific Reports}\\
Nature Publishing Group

\bigskip

Dear Editor,

I submit for your consideration the manuscript titled \textbf{``Predicting
Speech-in-Noise Performance Beyond the Pure-Tone Audiogram: A Machine
Learning Benchmark on the Oldenburg Hearing Health Record''} for publication
as an Article in \textit{Scientific Reports}.

\textbf{What this paper does.}
Hearing loss affects 1.6 billion people globally, yet its standard clinical
assessment---the pure-tone audiogram---explains only part of real-world
speech understanding.  I present what is, to our knowledge, the first
systematic machine learning benchmark on the recently published Oldenburg
Hearing Health Record (OHHR; Kollmeier et al., \textit{Scientific Data}
2025), an open CC~BY~4.0 dataset of 581 adults with audiograms, adaptive
loudness scaling, speech-in-noise tests, and cognitive and self-report
measures.  The study systematically evaluates five progressively richer
feature blocks (1--41 features) across Ridge regression, Random Forest,
LightGBM, and a multilayer perceptron under 10-fold cross-validation, and
includes SHAP-based feature importance, severity-stratified analysis,
calibrated probabilistic classification, a minimal-predictor analysis, and
partial cross-cohort validation on NHANES 2017--18 ($N=859$).

\textbf{Key findings.}
A single binaural pure-tone average captures 71\% of speech-in-noise
variance; the full battery achieves $R^2=0.81$ (Random Forest,
MAE\,=\,1.23\,dB~SNR).  Critically, the benefit of extended audiological
assessment over the audiogram alone is concentrated in \emph{mild} hearing
loss (26--40\,dB), where the audiogram explains only 9\% of variance and
the full battery achieves 37\%---exactly the patient group where
audiogram-only counselling is most uncertain.  Age and loudness-growth
slope are the most informative features beyond the audiogram, a finding
that partially replicates in independent NHANES data (AUROC: 0.747
$\to$ 0.801 when age is added).  Classification of clinically poor
speech-in-noise performers achieves AUROC\,=\,0.978 with well-calibrated
probability estimates.

\textbf{Why \textit{Scientific Reports}.}
This work is quantitative, reproducible, and clinically relevant.  The
dataset is publicly available (Zenodo, DOI 10.5281/zenodo.16919812) and
all analysis code is released at
\url{https://github.com/ryanlearnscs50/oldenburg} under the MIT licence.
The study analyses a pre-existing anonymised public dataset and requires
no additional ethics approval.

\textbf{Suggested reviewers.}
\begin{itemize}
  \item \textbf{Prof.\ Tobias Dau}, Hearing Systems Group, Technical
        University of Denmark (DTU), Denmark.
        (\href{mailto:tobias.dau@dtu.dk}{tobias.dau@dtu.dk})\\
        \textit{Expertise: computational auditory modelling, speech
        intelligibility prediction, hearing science.}

  \item \textbf{Prof.\ Michael Akeroyd}, Hearing Sciences, School of
        Medicine, University of Nottingham, UK.
        (\href{mailto:michael.akeroyd@nottingham.ac.uk}{michael.akeroyd@nottingham.ac.uk})\\
        \textit{Expertise: hearing epidemiology, large-scale audiological
        datasets, quantitative hearing research.}

  \item \textbf{Prof.\ Jan Wouters}, ExpORL, Department of Neurosciences,
        KU~Leuven, Belgium.
        (\href{mailto:jan.wouters@kuleuven.be}{jan.wouters@kuleuven.be})\\
        \textit{Expertise: audiological assessment, hearing aids, cochlear
        implants, and auditory signal processing.}
\end{itemize}

\textbf{Competing interests.}
The author declares no competing financial or non-financial interests.

\textbf{AI assistance.}
Portions of the code and manuscript were drafted with the assistance of an
AI language model (Claude, Anthropic).  All scientific content, numerical
results, and interpretations were verified by the author, consistent with
\textit{Scientific Reports}' authorship and AI-use policies.

I confirm that this manuscript has not been published elsewhere and is not
under consideration by any other journal.

Yours sincerely,

\bigskip
Ryan Yi Sheng Neo\\
Independent Researcher\\
\href{mailto:realryanneo@gmail.com}{realryanneo@gmail.com}\\
\url{https://github.com/ryanlearnscs50/oldenburg}

\end{document}
